\documentclass[11pt,a4paper]{article}
\usepackage{fontspec,amsmath,amssymb,amsthm,microtype}
\setmainfont{lmroman10-regular.otf}[BoldFont=lmroman10-bold.otf,ItalicFont=lmroman10-italic.otf,BoldItalicFont=lmroman10-bolditalic.otf]
\usepackage[margin=25mm]{geometry}
\usepackage[unicode,hidelinks]{hyperref}
\hypersetup{pdfauthor={LangevinLean Project}}
\setlength{\emergencystretch}{3em}
\setlength{\parskip}{0.4em}
\setlength{\parindent}{0pt}
\allowdisplaybreaks[2]
\raggedbottom
\title{\texorpdfstring{The nonsmooth single-\allowbreak{}pair energy estimate}{The nonsmooth single-pair energy estimate}}
\author{LangevinLean Project}
\date{}
\begin{document}\maketitle
Draft mathematical exposition; these notes may contain mistakes.


The energy inequality compares the true proximal output with any comparator while retaining the original nonsmooth potential difference at the input. An input subgradient controls the difference between input and proximal output. Completing a square produces a forcing term proportional to the square of the proximal scale, the correct quantity for summing particle energy estimates.




For convex \(f\), a true proximal minimizer \(p=P_\eta(x)\), positive \(\eta\), and a genuine input subgradient \(s\in\partial f(x)\), proximity controls both squared distance and the original nonsmooth energy.


\section{\texorpdfstring{Expand supports and complete a square}{Expand supports and complete a square}}


Let \(p=P_\lambda(y)\), \(s\in\partial g_2(y)\), and \(a\) be any comparator. Proximal optimality supplies the residual subgradient at \(p\). Combine its support inequality at \(a\) with the support inequality for \(s\) at \(p\), expand the squared distances, and use the nonnegative square \(\|(y-p)-\lambda s\|^2\). The result is v3 Lemma A.8:


\[
\|p-a\|^2\le\|y-a\|^2
 -2\lambda\bigl(g_2(y)-g_2(a)\bigr)+\lambda^2\|s\|^2.
\]


If \(\|s\|^2\le C\), replace the last term by \(\lambda^2C\). The proof works with any genuine subgradient satisfying that bound; it does not assume an energy estimate or a convergence statement. Taking the minimal section in C1 is one way to supply such a vector.





Write \(d=x-p\). The residual support at \(p\) implies \(2\eta[f(p)-f(a)]\le2\langle d,p-a\rangle\). The input support gives \(f(x)-f(p)\le\langle s,d\rangle\). Add these to the squared-distance expansion. The remaining scalar expression is \(2\eta\langle s,d\rangle-\|d\|^2\le\eta^2\|s\|^2\), obtained from \(\|d-\eta s\|^2\ge0\).


\section{\texorpdfstring{Instantiate the chosen point and bound the forcing}{Instantiate the chosen point and bound the forcing}}


The existence theorem proves that the chosen \(P_\eta(x)\) meets the variational premise. Substitution gives \(\|P_\eta(x)-a\|^2\le\|x-a\|^2-2\eta(f(x)-f(a))+\eta^2\|s\|^2\). If \(\|s\|^2\le C\), multiply by the nonnegative \(\eta^2\) to replace the final term by \(\eta^2 C\). The vector must satisfy the true support inequality; a merely named subgradient is insufficient.




Expand \(\|p-y\|^2=\|x-y\|^2-2\langle x-p,x-y\rangle+\|x-p\|^2\), and use the residual subgradient at \(p\) to replace its pairing with the potential difference. A supplied subgradient \(s_x\in\partial f(x)\) bounds the remaining difference \(f(x)-f(p)\) by \(\langle s_x,x-p\rangle\). Completing the square gives \(2\lambda\langle s_x,x-p\rangle-\|x-p\|^2\le\lambda^2\|s_x\|^2\). Thus a bound \(\|s_x\|^2\le C\) produces the forcing \(\lambda^2C\). This algebra is why that forcing has a squared proximal scale.





The comparator is arbitrary, so the same deterministic inequality can later be applied with a random reference point pointwise. A bound on \(\|s\|^2\) supplies only the last forcing term; it does not replace the requirement that \(s\) be a genuine supporting subgradient. The positive scale ensures both actual minimizer existence and validity of the residual optimality characterization used in the argument.


\end{document}
